% =====================================================================
%  2bat-ud1-3.tex  —  SESSIÓ 3: Llegir una matriu: dimensions i elements
%  (sense preàmbul: s'inclou des de main.tex)
%  Criteris: C1 (construir i llegir) i C4 (interpretar una taula com a matriu).
% =====================================================================
\horatitol{Sessió 3 — Llegir una matriu: dimensions i elements}%
{Consolidar la lectura d'una matriu: dimensió, posició d'un element i lectura per files i per columnes.}

\subsection{Idea clau}

A la sessió anterior vam construir quadres de comandament. Ara els donem nom
matemàtic i n'aprenem a \textbf{llegir} la informació amb precisió. Una matriu no
és només una taula bonica: cada nombre ocupa una \textbf{posició} concreta, i
saber-la indicar ens permet parlar de les dades sense embolics.

La clau és llegir en els \emph{dos sentits}: una \textbf{fila} ens explica una
mateixa cosa al llarg del temps o dels conceptes (per exemple, un servei mes a
mes), i una \textbf{columna} ens compara totes les coses en un mateix moment (per
exemple, tots els serveis en un mes concret).

\begin{idea}
\textbf{Vocabulari bàsic}
\begin{itemize}[leftmargin=1.4em, itemsep=2pt]
  \item \textbf{Matriu:} taula de nombres ordenats en \textbf{files} (horitzontals)
        i \textbf{columnes} (verticals). Sovint l'anomenem amb una lletra
        majúscula: $A$, $B$, $H$\dots
  \item \textbf{Dimensió:} s'escriu \emph{files $\times$ columnes}. Una matriu amb
        $3$ files i $4$ columnes té dimensió $3\times 4$ i, per tant, $12$ elements.
  \item \textbf{Element:} cada nombre de la matriu. Es localitza dient la
        \textbf{fila} i la \textbf{columna} on és (sempre fila primer, columna després).
  \item \textbf{Matrius especials:} si té una sola fila és una \textbf{matriu fila};
        si té una sola columna, una \textbf{matriu columna}; si té tantes files
        com columnes ($n\times n$), una \textbf{matriu quadrada}.
\end{itemize}
\end{idea}

\subsection{Exemples}

\begin{exemple}[llegir per files i per columnes]
Un centre obert anota els infants atesos en tres serveis (reforç, lleure i menjador)
durant tres mesos:
\[
A=\begin{pmatrix} 24 & 26 & 30 \\ 18 & 20 & 22 \\ 40 & 38 & 45 \end{pmatrix}.
\]
És una matriu de dimensió $3\times 3$ (i, per tant, quadrada). La \textbf{fila 1}
($24,\;26,\;30$) és el reforç al llarg dels tres mesos. La \textbf{columna 3}
($30,\;22,\;45$) són tots els serveis del tercer mes. L'element de la
\textbf{fila 2, columna 1} és $18$: el lleure el primer mes.
\end{exemple}

\begin{exemple}[matrius fila i columna]
Les places d'un únic servei durant quatre setmanes formen una \textbf{matriu fila}
de dimensió $1\times 4$:
\[
F=\begin{pmatrix} 12 & 15 & 14 & 18 \end{pmatrix}.
\]
Si en canvi anotem les hores de quatre professionals en un sol servei, tenim una
\textbf{matriu columna} de dimensió $4\times 1$:
\[
C=\begin{pmatrix} 10 \\ 8 \\ 12 \\ 6 \end{pmatrix}.
\]
\end{exemple}

\subsection{Exercicis resolts}

\begin{exresolt}[dimensió, element i lectura]
Una associació recull les persones ateses en dos punts d'acollida durant quatre
mesos:
\[
P=\begin{pmatrix} 50 & 55 & 60 & 58 \\ 30 & 28 & 35 & 40 \end{pmatrix}.
\]
La fila indica el punt d'acollida (Punt 1 i Punt 2) i la columna, el mes
(de gener a abril).
\begin{enumerate}[label=\alph*), leftmargin=1.6em, topsep=2pt]
  \item Quina dimensió té? És quadrada?
  \item Quin és l'element de la fila 1, columna 3? Què significa?
  \item Quantes persones va atendre el Punt 2 al mes d'abril?
\end{enumerate}
\solucio
a) Té $2$ files i $4$ columnes: dimensió $2\times 4$. No és quadrada (el nombre de
files i de columnes no coincideix).

b) La fila 1, columna 3 és $60$: les persones ateses al Punt 1 el mes de març.

c) El Punt 2 és la fila 2 i l'abril és la columna 4: l'element és $40$.
\resposta{Dimensió $2\times 4$ (no quadrada); \;fila 1, columna 3 $=60$ (Punt 1, març); \;Punt 2 a l'abril: $40$ persones.}
\end{exresolt}

\begin{exresolt}[treure conclusions d'una matriu]
Les hores setmanals de tres professionals (educadora, treballadora social i
monitora) en dos serveis són
\[
H=\begin{pmatrix} 14 & 6 \\ 9 & 11 \\ 5 & 13 \end{pmatrix}.
\]
La fila indica el professional (en aquest ordre) i la columna, el servei
(Servei 1 i Servei 2).
\begin{enumerate}[label=\alph*), leftmargin=1.6em, topsep=2pt]
  \item Indica'n la dimensió.
  \item Quantes hores fa en total la treballadora social?
  \item Quin servei concentra més hores en conjunt?
\end{enumerate}
\solucio
a) Té $3$ files i $2$ columnes: dimensió $3\times 2$.

b) La treballadora social és la fila 2: $9 + 11 = 20$ hores.

c) Sumem cada columna. Servei 1: $14 + 9 + 5 = 28$ hores. Servei 2:
$6 + 11 + 13 = 30$ hores. El Servei 2 concentra més hores ($30$ contra $28$).
\resposta{Dimensió $3\times 2$; \;treballadora social: $20$ hores; \;el Servei 2 concentra més hores ($30$ contra $28$).}
\end{exresolt}

\subsection{Exercicis per fer}

\begin{exercicis}
\begin{exlist}
  \item Considera la matriu
        \[ B=\begin{pmatrix} 8 & 5 & 3 \\ 6 & 9 & 7 \end{pmatrix}. \]
        \begin{enumerate}[label=\alph*), leftmargin=1.6em, topsep=2pt]
          \item Quina dimensió té? És quadrada?
          \item Quin és l'element de la fila 2, columna 1?
          \item Quin és l'element de la fila 1, columna 3?
        \end{enumerate}
  \item Un casal anota els infants inscrits en quatre tallers durant un trimestre:
        \[ T=\begin{pmatrix} 20 & 18 & 22 & 19 \end{pmatrix}. \]
        Quina dimensió té? Com s'anomena aquest tipus de matriu? Quants tallers hi
        ha en total?
  \item Una matriu recull les places ocupades de tres centres (files) en dos
        trimestres (columnes):
        \[ Q=\begin{pmatrix} 90 & 100 \\ 70 & 65 \\ 110 & 120 \end{pmatrix}. \]
        \begin{enumerate}[label=\alph*), leftmargin=1.6em, topsep=2pt]
          \item Indica'n la dimensió.
          \item Quin centre té més ocupació el segon trimestre?
          \item En quin trimestre, en conjunt, hi ha hagut més places ocupades?
        \end{enumerate}
  \item Una matriu té dimensió $5\times 3$. Quantes files té? Quantes columnes?
        Quants elements té en total? Podria ser quadrada?
  \item Escriu \emph{tu mateix\slash a} una matriu columna de dimensió $3\times 1$
        que representi les hores de voluntariat de tres persones en una setmana, i
        explica què vol dir cada element.
  \item Indica la dimensió de cadascuna d'aquestes matrius i digues de quin tipus
        és (fila, columna, quadrada o cap d'aquests):
        \[ M=\begin{pmatrix} 4 & 7 \\ 2 & 5 \end{pmatrix},\quad
           N=\begin{pmatrix} 3 \\ 9 \\ 1 \end{pmatrix},\quad
           R=\begin{pmatrix} 6 & 8 & 2 \end{pmatrix},\quad
           S=\begin{pmatrix} 1 & 0 & 5 \\ 4 & 3 & 2 \end{pmatrix}. \]
  \item \textbf{(Raonament)} La matriu
        \[ D=\begin{pmatrix} 35 & 40 & 38 \\ 50 & 48 & 55 \end{pmatrix} \]
        recull els àpats servits en dos menjadors socials (files) durant tres dies
        (columnes). Sense calcular-ho tot, explica com decidiries quin menjador ha
        servit més àpats en total i quin dia, en conjunt, se n'han servit més.
        Després comprova-ho amb els nombres.
\end{exlist}
\end{exercicis}

% ---------------------------------------------------------------------
\fullsolucions{Sessió 3}
\begin{solucions}
\begin{sollist}
  \item (a) dimensió $2\times 3$, no és quadrada \quad (b) $6$ \quad (c) $3$
  \item Dimensió $1\times 4$; és una \textbf{matriu fila}; hi ha $4$ tallers.
  \item (a) dimensió $3\times 2$ \quad (b) el tercer centre, amb $120$ \quad
        (c) el segon trimestre: $100+65+120=285$ contra $90+70+110=270$ del primer.
  \item $5$ files, $3$ columnes, $15$ elements. No pot ser quadrada (files i
        columnes són diferents).
  \item Resposta oberta. Exemple vàlid:
        $\begin{pmatrix} 4 \\ 6 \\ 3 \end{pmatrix}$, on cada element són les hores
        setmanals de voluntariat d'una persona.
  \item $M$: $2\times 2$, quadrada. \;$N$: $3\times 1$, matriu columna. \;
        $R$: $1\times 3$, matriu fila. \;$S$: $2\times 3$, cap d'aquests tipus.
  \item Per saber quin menjador ha servit més àpats se sumen les files
        ($35+40+38=113$ i $50+48+55=153$): el segon menjador, amb $153$. Per saber
        el dia amb més àpats se sumen les columnes ($85$, $88$ i $93$): el tercer
        dia, amb $93$.
\end{sollist}
\end{solucions}
